\chapter{Critical analysis}\label{ch:critical-analysis}

This chapter steps back from the mechanics of SR-FCA and asks what the paper is and is not saying. \S\ref{sec:strengths} lists the strengths -- what the method buys us over prior work. \S\ref{sec:weaknesses} lists the weaknesses, including some that only become visible once one tries to run the algorithm in practice. \S\ref{sec:assumptions} unpacks the implicit assumptions and discusses their realism for modern deep-learning workloads. \S\ref{sec:future} closes with future directions.

\section{Strengths}\label{sec:strengths}

\paragraph{S1 -- $K$ is an output, not an input.} The single most practically important property of SR-FCA is that the number of clusters is discovered. For every prior clustered-FL method, choosing $K$ requires either an oracle (which exists only in simulated experiments) or running the full pipeline multiple times and selecting on a held-out validation loss. The latter is expensive: an IFCA sweep over $K \in \{2, 3, 4, 5\}$ costs $4\times$ the compute of a single run. SR-FCA pays this price once, through the pairwise-distance histogram that sets $\lambda$ -- a single cheap pass at $T_0$ local steps -- and from then on the graph topology picks $K$ for us. Our experiments (Chapter~\ref{ch:experiments}) confirm that on both MNIST splits the histogram has a wide, clean valley between intra- and inter-cluster modes (Figure~\ref{fig:pairwise-hist}), so the valley-midpoint rule recovers the correct $K^\star$ ($4$ on rotated, $2$ on inverted) on every seed without per-seed tuning. The same observation transfers to Fashion-MNIST in Chapters~\ref{ch:experiments} and~\ref{ch:part-c}, where the histogram heuristic again succeeds on the first try.

\paragraph{S2 -- no initialisation assumption.} IFCA's theoretical guarantees require that the initial $K$ models already lie in the basin of attraction of the true cluster minimisers; this is achieved in practice by random restarts, which are wasteful and brittle. SR-FCA starts from a \emph{shared} random initialisation $w_0$ broadcast to all clients, and the ONE\_SHOT graph is built on the local models that result from a short local-SGD pass. Because every client departed from the same $w_0$, the pairwise distances reflect data heterogeneity rather than initialisation noise. This is visible in our pairwise-distance histograms (Figure~\ref{fig:pairwise-hist}): the intra-cluster and inter-cluster modes are cleanly separated on both MNIST-rotated and MNIST-inverted, and the valley between them is the appropriate $\lambda$.

\paragraph{S3 -- robustness to mis-clustered clients.} TrimmedMeanGD is robust up to a $\beta$ fraction of adversarially-corrupted clients per coordinate (Yin et al.~\parencite{yin2018trimmedmean}). In SR-FCA this robustness is repurposed to tolerate \emph{mis-clustered} clients: even if ONE\_SHOT assigns a client to the wrong cluster, its divergent update is trimmed away at aggregation, the cluster model stays close to $w_c^\star$, and the next RECLUSTER call moves the client to the correct cluster. This is why Theorem~4.14 guarantees exact cluster recovery after only one REFINE round. Empirically the ONE\_SHOT misclustering rate was zero on every seed of our simulated data -- both rotation and inversion splits, on both MNIST and Fashion-MNIST -- so the robustness margin that TrimmedMeanGD provides is never exercised in practice. It is theoretically necessary for the correctness proof but empirically dormant on idealised image-transform benchmarks.

\section{Weaknesses}\label{sec:weaknesses}

\paragraph{W1 -- $\mathcal{O}(m^2)$ pairwise-distance cost at ONE\_SHOT.} The graph construction requires computing $\binom{m}{2}$ pairwise distances, each of cost $\Theta(P)$ (flat parameter size). For $m \le 200$ clients -- the paper's regime -- this is trivial. For $m = 10^4$ clients with $P = 10^6$ parameters (realistic mobile-FL scale), pairwise $\ell_2$ alone is $5 \times 10^{13}$ operations, and the cross-cluster loss distance is worse because each pair involves a forward pass. IFCA scales with $K m$ rather than $m^2$. The paper does not discuss this; it becomes the main obstacle to applying SR-FCA at large scale.

\paragraph{W2 -- worse statistical rate than IFCA.} Theorem~4.8 gives $\tilde{\mathcal{O}}(1/\sqrt{n})$ per-cluster error, against IFCA's $\tilde{\mathcal{O}}(1/n)$ (Remark~4.16 in~\parencite{vardhan2024srfca}). This is a real cost in the limit of many samples per client: for large $n$, a well-initialised IFCA will produce strictly better per-cluster models. The practical consequence is that SR-FCA should be preferred in the \emph{cold-start} regime (no good initialisation available, $K$ unknown), while IFCA should be preferred once the clustering is established. The paper does not propose a hybrid, and we think this is an obvious missed opportunity (see F1 below).

\section{Implicit assumptions}\label{sec:assumptions}

The SR-FCA guarantees rely on three assumptions stated in the paper's \S 4 (Assumptions 4.3--4.5 in~\parencite{vardhan2024srfca}). We discuss each in turn.

\paragraph{Strong convexity (A4.3).} Each client's empirical risk $f_i$ is assumed $\mu$-strongly convex. This holds for ridge regression, logistic regression with $\ell_2$ regularisation, and strictly convex surrogates, but \emph{fails} for neural networks with ReLU activations and for almost every modern deep model. The paper addresses this pragmatically: for the MNIST/CIFAR experiments it substitutes $\ell_2$ distance with the cross-cluster loss distance~\eqref{eq:ccloss}, which is a meaningful distance even on non-convex landscapes. However the formal guarantees do not transfer. In our view this is fine -- clustered-FL is a cluster-recovery problem, and cluster recovery in practice uses the cross-cluster loss because two good models of the same task are \emph{functionally close} even if they sit at different minima of the same landscape.

\paragraph{Smoothness (A4.4).} $f_i$ has $L$-Lipschitz gradients. This is a benign assumption that fails only for architectures with non-differentiable operators at scale (e.g., hard attention). For our 2-layer MLP with ReLU it is effectively satisfied except on a measure-zero set, which is the standard excuse in the federated-optimisation literature.

\paragraph{Cluster separation (A4.5).} Each pair of ground-truth minimisers $(w^\star_c, w^\star_{c'})$ is at distance at least $\Delta > 0$, and $\Delta$ is bounded below by a function of $n$, $m$, $\mu$, $L$. This is essential -- without separation there is no clustering to recover. Crucially, the paper states the assumption on $\{w^\star_c\}$, i.e., on the ground-truth minimisers, whereas the algorithm only sees the empirical minimisers $\{\hat{w}_c\}$. Estimation error contracts $\Delta$ to an \emph{observed} separation $\hat{\Delta}$, and the paper's $\lambda$ calibration assumes $\hat{\Delta}$ is large enough. In our experiments the assumption is easy to validate post hoc by looking at the pairwise histogram; if the two modes do not separate, no choice of $\lambda$ will do.

\paragraph{Bounded heterogeneity.} Not listed explicitly but implicit: clients inside the same cluster are assumed IID. Our simulated data satisfy this by construction -- the rotated/inverted transformation is applied identically to every client's images inside the cluster. On real federated data (FEMNIST, Shakespeare) the assumption is only approximate, and the paper's experiments demonstrate empirically that SR-FCA is robust to modest intra-cluster heterogeneity. We do not evaluate this property -- it is outside the core scope -- and flag it as a direction for follow-up work.

\section{Future directions}\label{sec:future}

\paragraph{F1 -- hybrid SR-FCA$\to$IFCA.} W2 above already hints at this: use SR-FCA to discover $K^\star$ and a near-correct clustering, then hand off to IFCA to train the cluster models to the faster $\tilde{\mathcal{O}}(1/n)$ rate. The handoff is mechanically trivial (SR-FCA already produces the initial models IFCA needs); the analytical contribution would be to show that the $\tilde{\mathcal{O}}(1/n)$ rate is inherited without re-introducing IFCA's initialisation brittleness. The paper's own conclusion flags this possibility briefly but does not pursue it.

\paragraph{F2 -- FedProx aggregation inside REFINE.} The paper uses TrimmedMeanGD for the robust inner loop. FedProx~\parencite{li2020heterogeneity} achieves robustness via a proximal penalty on client updates rather than coordinate-wise trimming; it is known to be more effective when clients are \emph{heterogeneous in system capability} (e.g., some clients drop out part-way through a round). Plugging FedProx into SR-FCA would be a one-line change to \texttt{src/federated.py} in our implementation and is worth an empirical comparison.

\paragraph{F3 -- hierarchical / nested clustering.} The MERGE step is a flat merge: any two clusters whose models are within $\lambda$ become one. A hierarchical variant could maintain a dendrogram and let downstream applications choose the granularity. This is especially natural when the ground-truth heterogeneity has nested structure (e.g., country $\to$ region $\to$ city for a mobile-FL deployment).

\paragraph{F4 -- scalable ONE\_SHOT via sketching.} The $\mathcal{O}(m^2 P)$ cost of ONE\_SHOT (W1) can in principle be reduced to $\mathcal{O}(m P + m^2 s)$ for sketch dimension $s \ll P$, using Johnson--Lindenstrauss random projections. The concentration of $\ell_2$ distances under JL is a textbook result; the open question is whether the sketched distances still support the threshold-graph guarantees of Lemma~4.6.
